\documentclass[ ../main.tex]{subfiles}
\providecommand{\mainx}{..}
\begin{document}
\section{Primitive approximate \emph{values}}
\label{sec:prims}
The following are the primitive algebraic data types and operations on types.

These are \emph{first-order} approximate types since the approximation has a single error rate per element.

\subsection{Void type}
The most primitive type is the \emph{empty set} type, denoted by \Void.
There are no elements in the empty set and therefore it is not possible to construct values of this type.
There is only one function in the set $\Void \to \Set{X}$, known as the \emph{absurd} function since it can never be invoked.
Since \Void has no values, $\AT{\Void}$ is equivalent to \Void.

$\AT{\Void}$ is necessary to complete the algebra of algebraic data types, but serves only a theoretical purpose.

\subsection{Unit type}
The \emph{unit} type, denoted by \Unit, is isomorphic to any set with only a single element, i.e., any \emph{singleton set}.
The set $\Unit \to \Set{X}$ has a cardinality $\Card{\Set{X}}$ and the set $\Set{X} \to \Unit$ has a cardinality $1$.
If we are talking about \emph{partial} functions, then $\Set{X} \pfun \Unit$ has cardinality $2^{\Set{X}}$.

The set $\Unit \to \Unit$ has a cardinality of $1$, which is the identity function $\IdFun \colon \Unit \to \Unit$.

The approximate unit type $\AT{\Unit}$ is necessary to complete the algebra, but given the \Unit type, similar to the \Void type there is no uncertainty about its value, i.e., $\AT{\Unit}$ is equivalent to \Unit.

In addition, if there is some function $\Fun{f} \colon \Unit \to \Set{X}$ its approximate analog is $\APFun{f} \colon \Unit \to \Set{X}$, which models an \emph{approximate} constant.

\subsection{Composite type constructors}

Sum, product, and exponential types are composite constructors treated in \cref{sec:composite}.

\begin{remark}[Partial functions]
Partial functions (functions not defined on their entire domain) introduce an additional source of approximation error: elements outside the preimage may be mapped to some element in the codomain rather than correctly yielding no result.
The companion paper on Bernoulli maps~\cite{bernoulliMaps} develops the full theory of approximate partial functions, including false mapping rates and their information-theoretic lower bounds on space complexity.
\end{remark}

\end{document}